% =====================================================================
%  abstract.tex  --  PHASE VII, Task 14E
%
%  Rewritten from the Phase I draft.  Four things changed and each is a
%  consequence of a measured result rather than a stylistic choice:
%
%   1. The forecasting component is a gradient-boosted tree, not the
%      Seq2Seq LSTM.  Phase V: the sequence model is beaten by an
%      untrained persistence reference by 7.95 s of per-window MAE.
%   2. The decision-layer saving is 2.9 +/- 0.6 USD/yr, not 7.2, and it
%      is NOT claimed as a saving -- its interval spans zero.  The claim
%      is the paired contrast against the static rule.
%   3. The calendar share of the duration forecaster is 82% (Shapley),
%      not 94.5% (split gain).
%   4. The physics accept/reject loop is named as a contribution.
%
%  NO placeholder survives.  Every number below is traceable to a
%  stored output; see the section named beside it in the source.
% =====================================================================

\begin{abstract}
Industrial machines consume a substantial fraction of their electrical
energy while energised but not producing, yet eliminating this standby
consumption in an existing facility is obstructed by three coupled
difficulties: per-timestamp operating-state annotations do not exist and
are prohibitively expensive to obtain; the duration of a future
non-productive interval is unknown at the moment a shutdown decision must
be made; and shutting a machine down is economical only when that
duration exceeds a break-even point set by restart cost and standby
power. Existing work addresses these in isolation---non-intrusive load
monitoring infers states but neither forecasts nor acts, while
energy-aware scheduling decides optimally but assumes the idle duration
is given. This paper proposes a unified, label-free framework that closes
the loop from raw electrical measurements to an economically justified
shutdown decision, and evaluates it end to end on IMDELD, a public
benchmark of $1$\,Hz measurements from eight heavy machines in a
full-scale plant. The framework comprises (i)~\emph{physics-gated
unsupervised state inference}, in which a Gaussian hidden Markov model is
fitted without labels and the resulting partition is then \emph{accepted
or rejected} against circuit-level checks external to the model;
(ii)~\emph{non-productive duration forecasting} benchmarked against an
untrained reference; and (iii)~\emph{optimisation-based decision
support}, in which the break-even duration is derived as the solution of
a constrained cost minimisation rather than supplied as a threshold. The
accept/reject step is load-bearing: fitted at ten seeds the labeller
places the standby cluster on a physically different state in three, and
conditioning on a passing score collapses the standby-hour spread from
$\pm33.7\%$ to $\pm3.0\%$. Temporal decoding reduces state flicker from
$66.9\%$ to $44.5\%$ and an explicit minimum-dwell constraint removes it
entirely, at a cost of $0.19\%$ of rows relabelled. Two results are
reported against expectation: of seven forecasters, five neural sequence
models are beaten by an untrained persistence reference, and only a
gradient-boosted tree clears it ($11.50$ against $12.95$\,s of
per-window MAE, $p_{\text{Holm}}=7\times10^{-4}$); and the proposed
policy cannot be shown to save money against the plant's own behaviour on
an 18-day held-out sample ($+2.9$\,USD/yr, $95\%$ CI $[-56,+46]$), though
it beats the fixed-threshold rule it replaces by $+68.9$\,USD/yr
($p=0.001$). Across the plant the state definition is physically valid on
the four motor-driven machines and correctly rejected on the four
auxiliary ones, and the facility's total attainable standby saving is
under $90$\,USD/yr. The framework's demonstrated value on this facility
is therefore decision quality and validated state inference rather than
recovered energy, and it requires only the voltage and current channels a
standard retrofit meter already provides.
\end{abstract}
